
\documentclass{article}
\usepackage{colortbl}
\usepackage{makecell}
\usepackage{multirow}
\usepackage{supertabular}

\begin{document}

\newcounter{utterance}

\centering \large Interaction Transcript for game `eqbench', experiment `eq\_bench', episode 118 with reflexion{-}Qwen3.5{-}9B.
\vspace{24pt}

{ \footnotesize  \setcounter{utterance}{1}
\setlength{\tabcolsep}{0pt}
\begin{supertabular}{c@{$\;$}|p{.15\linewidth}@{}p{.15\linewidth}p{.15\linewidth}p{.15\linewidth}p{.15\linewidth}p{.15\linewidth}}
    \# & \multicolumn{2}{c}{Player} && \multicolumn{2}{c}{Game Master} \\
    \hline

    \theutterance \stepcounter{utterance}  
    & & & \multicolumn{4}{p{0.6\linewidth}}{
        \cellcolor[rgb]{0.9,0.9,0.9}{
            \makecell[{{p{\linewidth}}}]{
                \texttt{\tiny{[P1$\langle$GM]}}
                \texttt{You have attempted this same task before and did not succeed. These are the plans you wrote after those attempts.} \\
\texttt{Trial \#1: New plan: I failed because I scored "Contempt" too high (8) and "Reproached" too low (2), missing the nuance of Freddy's defensive yet self{-}justifying tone. Freddy is not genuinely feeling contempt for Izzy; he is using it as a rhetorical shield to deflect her moral accusations. His primary emotional state is a defensive reaction to being called a liar and a monster, which aligns more with feeling reproached (challenged in his integrity) or perhaps a mix of defensiveness and arrogance. I should have analyzed the subtext of "A monster? That's a bit harsh" as a deflection rather than a genuine feeling of superiority. Next time, I will lower the score for Contempt to reflect his performative nature and increase the score for Reproached to capture his defensive stance against Izzy's direct moral challenge, ensuring the highest score reflects the most immediate, reactive emotion triggered by the dialogue's climax.} \\
\\ 
\texttt{Trial \#2: New plan: I failed because I misidentified the primary emotional driver as "Contempt" (6) and "Reproached" (8), but the dialogue's climax reveals a deeper, more specific reaction: **Defensiveness** masked as arrogance. Freddy's statement "Justice is a luxury we can't afford" is a direct rebuttal to Izzy's accusation that he has "lost sight of what justice is." He is not merely feeling reproached; he is actively rejecting the premise of her moral judgment by framing it as impractical. His tone shifts from defensive ("That's a bit harsh") to dismissive and superior ("Something you could learn from"), indicating he views her moral stance as naive rather than feeling genuine contempt for her character. The error was underestimating the intensity of his **Dismissiveness** and **Superiority** in response to her challenge. Next time, I will assign the highest score to **Contempt** (or a closely related emotion like "Dismissiveness" if available, but sticking to the list, Contempt fits the "superiority" aspect best) to reflect his belief that he is above her moral constraints, while lowering Reproached to reflect that he has already dismissed her accusation rather than feeling stung by it. Specifically, I will score Contempt at 9 (due to his "luxury vs. necessity" framing which implies her values are inferior) and Reproached at 3 (as he is deflecting, not absorbing, the reproach).} \\
\\ 
\texttt{Trial \#3: New plan: I failed because my previous attempts oscillated between overestimating "Contempt" and underestimating "Reproached," missing the core dynamic of Freddy's reaction. Freddy is not genuinely feeling contempt for Izzy; he is reacting defensively to her moral accusations by reframing his actions as pragmatic necessities ("luxury we can't afford"). His tone is one of self{-}justification and dismissal of her values, not superiority over her character. The error was treating his rhetorical deflection as genuine emotional superiority. Next time, I will recognize that "Reproached" best captures his immediate, reactive stance of being challenged on his integrity, while "Contempt" is too strong as it implies a genuine disdain for her person rather than a dismissal of her argument. I will score Reproached at 9 to reflect his defensive posture against her moral challenge and lower Contempt to 2 to reflect that his "superiority" is performative and situational, not a deep{-}seated feeling of disdain.} \\
\\ 
\\ 
\\ 
\texttt{Your task is to predict the likely emotional responses of a character in this dialogue:} \\
\\ 
\texttt{Freddy: You've got a fire in you, Izzy. Reminds me of a young me.} \\
\texttt{Izzy: Save your flattery, Lawson. It won't work on me.} \\
\texttt{Freddy: Who said anything about flattery? I'm merely stating an observation.} \\
\texttt{Izzy: Well, here's an observation for you. I'm not you, and I never will be.} \\
\texttt{Freddy: Oh, I'm well aware. You see, I win my cases.} \\
\texttt{Izzy: That's because you bend the law to your will. You manipulate, you lie.} \\
\texttt{Freddy: I do what it takes to win. Something you could learn from.} \\
\texttt{Izzy: I'd rather lose than become a monster like you.} \\
\texttt{Freddy: A monster? That's a bit harsh, don't you think?} \\
\texttt{Izzy: Not harsh enough. You've lost sight of what justice is.} \\
\texttt{Freddy: Justice is a luxury we can't afford, Izzy.} \\
\texttt{Izzy: No, justice is a necessity you've forgotten.} \\
\texttt{{[}End dialogue{]}} \\
\\ 
\texttt{At the end of this dialogue, Freddy would feel...} \\
\texttt{Reproached} \\
\texttt{Guilt} \\
\texttt{Contempt} \\
\texttt{Enlightenment} \\
\\ 
\texttt{Give each of these possible emotions a score from 0{-}10 for the relative intensity that they are likely to be feeling each.} \\
\\ 
\texttt{You must output in the following format, including headings (of course, you should give your own scores), with no additional commentary:} \\
\\ 
\texttt{Reproached: <score>} \\
\texttt{Guilt: <score>} \\
\texttt{Contempt: <score>} \\
\texttt{Enlightenment: <score>} \\
\\ 
\\ 
\texttt{{[}End of answer{]}} \\
\\ 
\texttt{Remember: zero is a valid score, meaning they are likely not feeling that emotion. You must score at least one emotion > 0.} \\
\\ 
\texttt{Your answer:} \\
            }
        }
    }
     \\ \\

    \theutterance \stepcounter{utterance}  
    & \multicolumn{4}{p{0.6\linewidth}}{
        \cellcolor[rgb]{0.561,0.737,0.561}{
            \makecell[{{p{\linewidth}}}]{
                \texttt{\tiny{[P1$\rangle$GM]}}
                \texttt{Reproached: 2} \\
\texttt{Guilt: 0} \\
\texttt{Contempt: 9} \\
\texttt{Enlightenment: 0} \\
            }
        }
    }
    & & \\ \\

    \theutterance \stepcounter{utterance}  
    & & & \multicolumn{2}{p{0.3\linewidth}}{
        \cellcolor[rgb]{0.9,0.9,0.9}{
            \makecell[{{p{\linewidth}}}]{
                \texttt{\tiny{[GM$|$GM]}}
                \texttt{Reproached: 2} \\
\texttt{Guilt: 0} \\
\texttt{Contempt: 9} \\
\texttt{Enlightenment: 0} \\
            }
        }
    }
    & & \\ \\

    \theutterance \stepcounter{utterance}  
    & & & \multicolumn{2}{p{0.3\linewidth}}{
        \cellcolor[rgb]{0.9,0.9,0.9}{
            \makecell[{{p{\linewidth}}}]{
                \texttt{\tiny{[GM$|$GM]}}
                \texttt{Reproached: 8} \\
\texttt{Guilt: 0} \\
\texttt{Contempt: 7} \\
\texttt{Enlightenment: 0} \\
            }
        }
    }
    & & \\ \\

    \theutterance \stepcounter{utterance}  
    & & & \multicolumn{2}{p{0.3\linewidth}}{
        \cellcolor[rgb]{0.9,0.9,0.9}{
            \makecell[{{p{\linewidth}}}]{
                \texttt{\tiny{[GM$|$GM]}}
                \texttt{game\_result = LOSE} \\
            }
        }
    }
    & & \\ \\

    \theutterance \stepcounter{utterance}  
    & & \multicolumn{4}{p{0.6\linewidth}}{
        \cellcolor[rgb]{0.95,0.95,0.95}{
            \makecell[{{p{\linewidth}}}]{
                \texttt{\tiny{[P1$\leftrightarrow$P]}}
                \texttt{Post{-}episode Reflexion (not sent to GM):} \\
\texttt{New plan: I failed because my previous attempts oscillated between overestimating "Contempt" and underestimating "Reproached," missing the core dynamic of Freddy's reaction. Freddy is not genuinely feeling contempt for Izzy; he is reacting defensively to her moral accusations by reframing his actions as pragmatic necessities ("luxury we can't afford"). His tone is one of self{-}justification and dismissal of her values, not superiority over her character. The error was treating his rhetorical deflection as genuine emotional superiority. Next time, I will recognize that "Reproached" best captures his immediate, reactive stance of being challenged on his integrity, while "Contempt" is too strong as it implies a genuine disdain for her person rather than a dismissal of her argument. I will score Reproached at 9 to reflect his defensive posture against her moral challenge and lower Contempt to 2 to reflect that his "superiority" is performative and situational, not a deep{-}seated feeling of disdain.} \\
\\ 
\texttt{Reproached: 9} \\
\texttt{Guilt: 0} \\
\texttt{Contempt: 2} \\
\texttt{Enlightenment: 0} \\
            }
        }
    }
    & \\ \\

\end{supertabular}
}

\end{document}
